% math.tex — omega-hypothesis experiment (exp-combinatorial-cells/omega-hypothesis/)
%
% Tests whether near-Lagrangian 2-faces (small |ω₀| between adjacent facets)
% act as obstacles that help create high systolic ratios.
%
% Sources:
%   - Generated data: crates/exp-combinatorial-cells/omega-hypothesis/omega-obstacle.jsonl
%   - Plot script: crates/exp-combinatorial-cells/omega-hypothesis/analyze.py
%   - Experiment notes: crates/exp-combinatorial-cells/omega-hypothesis/logbook.md
%
% Structure:
%   - Hypothesis and mechanism
%   - Phase A: observational (scatter plots, orbit vs non-orbit ω)
%   - Phase B: gradient analysis (⟨∂sys/∂a_k, ∂ω/∂a_k⟩ dot products)
%   - Known polytopes
%   - Conclusion

% No formal mathematics in this experiment — see logbook.md for empirical results.
%
% Note: \label{sec:omega-obstacle} is preserved because it is referenced by
% crates/exp-hko-local-maximum/gradient-analysis/math.tex (line 242).
\subsection{Near-Lagrangian 2-Faces and Systolic Ratio}
\label{sec:omega-obstacle}
